% Replace a deliberately schematic OTP(circuit) with the real-copy protocol.
\begin{frame}{Encrypted MBQC}
  \centering
  \setbeamercolor{itemize item}{fg=CVQCTeal}
  \begin{tikzpicture}
    \path[use as bounding box] (-6.6,-3.15) rectangle (6.6,3.05);

    \node[anchor=north west,inner sep=0pt,text=CVQCInk] at (-6.0,2.65) {%
      \begin{minipage}{12cm}
        \begin{itemize}
          \item Need $\text{trap}\approx\text{real}$ $\Longrightarrow$ encrypt everything.
        \end{itemize}
      \end{minipage}%
    };

    \only<1>{
      % Qcircuit's stick labels need explicit space inside the delimiters.
      \node[inner sep=0pt,text=CVQCInk] at (0,0.10) {%
        \resizebox{12.5cm}{!}{%
          $\operatorname{OTP}\!\left(
            \vcenter{\hbox{%
              \raisebox{0pt}[\dimexpr\height+3mm\relax][\dimexpr\depth+3mm\relax]{%
                \hspace*{0.7cm}%
                $\Qcircuit @C=1em @R=1em {
                  \lstick{\ket{\psi}} & \ctrl{1} & \gate{HT^c} & \meter & \rstick{m}\cw \\
                  \lstick{\ket{+}} & \control\qw & \qw & \qw & \rstick{X^mHT^c\ket{\psi}}\qw
                }$\hspace*{2.3cm}%
              }%
            }}
          \right)$%
        }%
      };
    }

    \only<2>{
      \node[inner sep=0pt,text=CVQCInk] at (0,0.10) {%
        \resizebox{12.5cm}{!}{%
          \hspace*{2.1cm}%
          $\Qcircuit @C=1em @R=1em {
            \lstick{X^{x_1}T^{t_1}\ket{\psi}} & \ctrl{1} & \gate{HT^\delta} & \meter
              & \cgate{\oplus r} & \rstick{\beta}\cw \\
            \lstick{X^{x_2}T^{t_2}\ket{+}} & \control\qw & \qw & \qw
              & \qw & \rstick{X^{x'}T^{t'}HT^c\ket{\psi}}\qw
          }$\hspace*{2.6cm}%
        }%
      };
      \node[font=\large,text=CVQCInk] at (0,-1.80)
        {$\delta=(-1)^{x_1}(c-t_1)+4x_2+4r\pmod{8}$};
      \node[font=\large,text=CVQCInk] at (0,-2.55)
        {$x'=x_2\oplus\beta,\qquad
          t'=(-1)^\beta(t_2+4x_1)\pmod{8}$};
    }
  \end{tikzpicture}

  \note<1>{
    \footnotesize
    The prover must not distinguish trap copies from real copies. The
    approximation symbol refers to its joint view of inputs and messages,
    not to the unencrypted logical states. Hide
    the input states and the measurement instructions, including whether
    an instance is a test. OTP around the entire circuit is deliberately
    schematic shorthand for this goal, not a literal operation on a
    circuit or a claim that applying an input pad alone suffices. The
    circuit inside the parentheses is the plaintext real-copy example.
    The next build replaces this shorthand with its actual implementation.
  }
  \note<2>{
    \footnotesize
    This is the manuscript's Ciphertext Teleportation protocol. Each input
    has its own pad, with independent uniform x bits and t exponents in
    Z8. Commuting these pads through CZ keeps both x bits and changes the
    t coordinates to $t_1'=t_1+4x_2$ and $t_2'=t_2+4x_1$ modulo 8.
    The paper sends $\delta=(-1)^{x_1}(c-t_1'+4r)$ for a fresh random bit r.
    Expanding $t_1'$ gives the displayed formula: a signed multiple of 4
    equals the unsigned multiple modulo 8. The prover applies
    $HT^\delta$ and reports the raw measurement bit b. The classical
    XOR gate is the verifier's unmasking $\beta=b\oplus r$.
    The lower label keeps the standard OTP form. The formulas below the
    circuit show how its final keys depend on the decoded bit beta.
    Absorbing the byproduct $X^\beta$ into the CZ-updated pad gives
    $x'=x_2\oplus\beta$ and $t'=(-1)^\beta(t_2+4x_1)$ modulo 8,
    up to global phase.
    These are honest-output formulas. Mask trap copies analogously.
    Next: iterative Pauli twirling.
    Source: CVQC-paper/4\_Real\_Copy.tex, Ciphertext Teleportation protocol
    and Uniformity Preserving lemma.
  }
\end{frame}
